% 11 November 2026. Paper-only: exact selected Hartree physical residual
% source/jet identities at moving interfaces; original count 11/13.
\medskip
\noindent\textbf{Physical residual traces determine the moving interface jets.}
Let \(W\) be the full piecewise radial physical field, matched in \(C^1_r\)
at a moving interface \(r=L(s)>0\). For one-sided traces write
\([f]=f(s,L(s)+)-f(s,L(s)-)\). The full genuine inverse-square quartic
Hartree operator and its residual are
\[
 \mathcal N_{\HH}(W)=\kappa_{\HH}
   (|\cdot|^{-2}*|W|^4)|W|^2W,\qquad
 \mathscr R_W=iW_s+W_{rr}+\frac2rW_r+\mathcal N_{\HH}(W).
\]
The nonlinear convolution is evaluated on the \emph{same global field},
including all four-region cross interactions, not separately on
individual artificial pieces. For smooth, decaying physical packets the
potential and its radial derivative are continuous at the interface.
Differentiating the two matched moving traces gives
\[
 [W]=[W_r]=[W_s]=0,\qquad [W_{sr}]=-L'[W_{rr}].
\]
Define \(d_j=W_{r^j}(s,L-)-W_{r^j}(s,L+)\) for \(j=2,3\). Directly taking
the jumps of the true PDE residual and its radial derivative gives
\begin{equation}
 \boxed{\,d_2=-[\mathscr R_W],\qquad
 d_3=-[\partial_r\mathscr R_W]+(iL'-2/L)d_2.\,}
 \label{eq:hbl-selected-physical-interface-jets}
\end{equation}
Indeed \([\mathscr R_W]=[W_{rr}]\), whereas
\[
 [\partial_r\mathscr R_W]
   =i[W_{sr}]+[W_{rrr}]+\frac2L[W_{rr}]
   =-d_3-\left(\frac2L-iL'\right)d_2.
\]
The total time derivatives along the moving interface are
\begin{align*}
 d_2'&=-[\partial_s\mathscr R_W]-L'[\partial_r\mathscr R_W],\\
 d_3'&=-[\partial_s\partial_r\mathscr R_W]
       -L'[\partial_r^2\mathscr R_W]
       +(iL''+2L'/L^2)d_2+(iL'-2/L)d_2',
\end{align*}
assuming the corresponding classical one-sided higher traces.
These are \emph{identities}, not positive-time-order bounds.
If both adjoining restrictions are exact solutions of the same
complete physical Hartree PDE in a neighborhood, then
\(\mathscr R_W=0\), hence \(d_2=d_3=0\), so a \(C^3\) shell is unnecessary
at that interface. A \(C^1\) join and the selected two real speeds
alone do not establish the required residual trace vanishing.

If the genuine residual jumps do not vanish, use the existing Hermite
correction
\[
 P=h^2d_2\phi_2((r-L)/h)+h^3d_3\phi_3((r-L)/h)
 \quad (L<r<L+h),
\]
extended by zero, with the previously computed endpoint jets.
Then \(U=W+P\) satisfies
\[
 [U]=[U_r]=[U_{rr}]=[U_{rrr}]=[U_s]=[U_{sr}]=0,
 \qquad
 \boxed{[\mathscr R_U]=[\partial_r\mathscr R_U]=0.}
\]
The full \(\mathcal N_{\HH}(W+P)-\mathcal N_{\HH}(W)\) is nonlocal and
need not vanish at the patch endpoints: its \emph{jump} and the jump of
its first radial derivative vanish. This is a trace regularity
statement, not the strong \(s^{-9/4}\) residual gain, its time
derivative, the two-parameter Lipschitz estimate, or the dynamical
blowup law.

The next source-producing computation is therefore the one-sided
\([\mathscr R_W]\) and \([\partial_r\mathscr R_W]\) on each of the three
actual interfaces, retaining the original unique real speed selection.
Applying the identities to already \emph{assumed} small residual traces
would be circular. The algebraic four-piece nonlinear PDE telescope
(CI 42958494409) does not supply the missing decay.
This paragraph is paper-only and does not change the retained
original-environment acceptance \(11/13\).
